\chapter{Quantum Groups: Foundations}
\label{ch:quantum-groups}

This chapter collects the quantum group background required for Part~III and the CY landscape.  The quantized enveloping algebra $U_q(\frakg)$ and its $R$-matrix are the targets of the CY-to-chiral functor at the representation-theoretic level: Conjecture~CY-C predicts that the braided monoidal category $\Rep_q(\frakg)$ arises as $\Rep^{E_2}(A_\cC)$ for a CY category $\cC(\frakg, q)$.  The Drinfeld--Jimbo and FRT presentations provide distinct access points, corresponding respectively to the Koszul dual (bar cohomology generators and relations) and the ordered bar (RTT $R$-matrix formalism) of the chiral algebra.

\section{Quantized enveloping algebras}
\label{sec:quantized-enveloping}

This section reviews the Drinfeld--Jimbo quantized enveloping algebra $U_q(\frakg)$ and the Lusztig divided-power form $\dot{U}_q(\frakg)$, emphasising the features relevant to the CY programme: the triangular decomposition (which matches the bar filtration), the Hopf algebra structure (which matches the bar coproduct), and the specialisation at roots of unity (which connects to the Kazhdan--Lusztig equivalence of Chapter~\ref{ch:quantum-group-reps}).

\section{The $R$-matrix and Yang--Baxter equation}
\label{sec:r-matrix-ybe}

The universal $R$-matrix $\cR \in U_q(\frakg)^{\hat{\otimes}2}$ is the quantum group avatar of the collision $r$-matrix $r(z) = \Res^{\mathrm{coll}}_{0,2}(\Theta_A)$ from Volume~I.  This section develops the Yang--Baxter equation, quasi-triangular Hopf algebras, and the passage from the formal $R$-matrix to the trigonometric and elliptic solutions that appear in the CY landscape.  The key structural point: the $R$-matrix lives in arity~$(1,1)$ of the ordered bar $B^{\mathrm{ord}}(A)$, and its Baxterisation $R(z)$ is the $z$-dependent lift (Volume~II, DK bridge).

\section{Representation categories}
\label{sec:qg-rep-categories}

This section establishes that $\Rep_q(\frakg)$ is a braided monoidal category (ribbon category when $q$ is generic), with the braiding provided by the universal $R$-matrix.  At roots of unity, the category acquires a finite tensor subcategory (the Andersen--Jantzen--Soergel quotient) that is modular in the sense of Turaev.  The target: the CY category $\cC(\frakg, q)$ should produce this braided category as $\Rep^{E_2}(\Phi(\cC))$.

\section{Drinfeld--Jimbo vs Faddeev--Reshetikhin--Takhtajan}
\label{sec:dj-vs-frt}

The Drinfeld--Jimbo presentation uses generators $E_i, F_i, K_i$ and Serre relations; the FRT presentation uses the $R$-matrix equation $R T_1 T_2 = T_2 T_1 R$ directly.  From the bar-complex perspective (Volume~II, Chapter on Yangians): the DJ presentation corresponds to the Koszul dual algebra $A^!$ (generators = bar $H^1$, relations = bar $H^2$), while the FRT/RTT presentation corresponds to the ordered bar with its $R$-matrix in bidegree $(1,1)$.  The equivalence of the two presentations is a theorem of Faddeev--Reshetikhin--Takhtajan for finite type; the extension to affine and toroidal types is developed in the toric CoHA chapter (Chapter~\ref{ch:toric-coha}).
